\documentclass[10pt,a4paper]{article}
\usepackage[margin=22mm]{geometry}
\usepackage{amsmath,amssymb,booktabs,array}
\usepackage[colorlinks=true,linkcolor=blue,urlcolor=blue]{hyperref}
\usepackage{microtype}
\newcommand{\R}{\mathbb{R}}
\title{RM9: Compact local merge state and strong-observer audit}
\author{Method-first ramp-metering review}
\date{1 October 2026}
\begin{document}
\maketitle
\noindent\textbf{Status.} RM8 Candidate A was approved. This is the first bounded
RM9 analytic audit. It derives conditional constraints, not an adopted simulator,
a full published-scheme validation, or a useful-gradient certificate.
A new finite-receiving allocation decision is pending.

\section{Contracts and proof ownership}
The task remains one metered finite-storage ramp entering a local multilane
mainline. Compare common bounded open-loop rates and count all affected traffic
through common clearance. The aggregate observer has the declared road-total
density, flow, speed and count histories, ramp histories, known lane parameters,
geometry, forcing and prior rates. Observation functions and source equations
must be specified before making an exact-history claim.

The conservation structure below is consistent with the lane-model lineage [1].
Projection, identifiability and accounting results are \emph{reviewer derivations}
under their stated conditions. They are not reported results of [1--4].
The full update in [1] includes additional modified-demand/supply rules;
none is silently repaired or certified here.

\section{A conserved sum need not be a closed state}
Let lane 1 be the shoulder lane; lanes 2 through $m$ are interior lanes.
For a cell of length $L$, $n_\ell=Lk_\ell$ is mass in vehicle equivalents,
$q^{\rm in/out}_\ell$ and $\lambda_{a\to b}$ are rates in vehicles per unit time.
Writing a generic conservative local balance, with realized ramp entry $r$,
\begin{equation}
 \dot n_\ell=q^{\rm in}_\ell-q^{\rm out}_\ell
 +\sum_a\lambda_{a\to\ell}-\sum_b\lambda_{\ell\to b}
 +\boldsymbol 1_{\ell=1}r. \label{eq:balance}
\end{equation}
Define shoulder/rest state $z=(n_1,n_R)$, with $n_R=\sum_{\ell=2}^m n_\ell$.
Internal transfers among rest lanes cancel. But the boundary transfer
$\lambda_{1\to2}$ depends on lane 2, not generally on $n_R$ alone.
For an autonomous exact projected update $G$, a necessary condition is
\begin{equation}
 P x=P x'\ \Longrightarrow\ P F(x,u)=P F(x',u)
 \quad\text{for every admissible common input }u. \label{eq:lump}
\end{equation}
Here $F$ is the detailed one-step map and $P$ the proposed projection.
If longitudinal/lateral source dependencies, mode memory or receiving states
are missing from $z$, summing balances does not prove this condition.

\paragraph{Explicit local-operator obstruction.}
For three or more lanes, compare spatially constant profiles in a local plateau
\begin{equation}
 (k_1,k_2,k_3)=(a,b,c),\qquad
 (k'_1,k'_2,k'_3)=(a,b+\epsilon,c-\epsilon),\quad \epsilon\ne0.
\end{equation}
Both have the same shoulder/rest masses. On the free-flow branch where the
source keep-right incentive is $I=1-b/a$, the weighted downstream densities
reduce to the constant local densities and the donor demand is $v_1a$.
Choose positive densities with $b,b+\epsilon<a/2$, all sending/receiving
limitations inactive. The candidate lateral-demand expression from [1] then
reduces to
\begin{equation}
 \lambda_{1\to2}(b)=v_1a\frac{a-2b}{a+b},\qquad
 \frac{d\lambda_{1\to2}}{db}=-\frac{3v_1a^2}{(a+b)^2}\ne0. \label{eq:lateral}
\end{equation}
This is an algebraic source-branch calculation, not a new transfer law.
Reverse transfer is inactive on this branch. A homogeneous plateau makes each
lane's longitudinal inflow/outflow balance zero locally; rest-internal
transfers cancel in $n_R$. Therefore the projected lateral evolution differs
although $z$ agrees. This rules out a \emph{general autonomous shoulder/rest
closure} on that branch without extra information or restrictions.
It does \textbf{not} rule out a closure conditioned on additional aggregate
observations, prove insufficiency of every scalar state, or establish identical
observable traffic histories. The two profiles above may have different road
flows and speeds.

\paragraph{Two distinct remedies.}
If the mainline has exactly two physical lanes, ``rest'' is one lane:
there is no within-rest aggregation. With known total $N=n_1+n_2$,
$n_1$ (or a contrast $n_1-n_2$) specifies the same lane information.
For three or more lanes, retain sufficient receiving-adjacent information,
prove a lossless reconstruction from admitted observables, or explicitly
adopt and justify an approximate mixing closure. The last option changes
the scientific claim and is not approved.

\section{Hidden at a snapshot versus hidden from the observer}
For a diagnostic equilibrium free-flow observation map, assume each lane has
known constant speed $v_\ell$, common observation support, no active supply
limit, and $q_\ell=v_\ell k_\ell$. Define road-total density $K$ and rate $Q$:
\begin{equation}
 K=\sum_\ell k_\ell,\qquad Q=\sum_\ell v_\ell k_\ell. \label{eq:moments}
\end{equation}
This condition is a diagnostic subcase, not a declaration that the entire
source update always realizes this observation map.
For two distinct lane speeds, already $K,Q$ determine the split:
\begin{equation}
 k_1=\frac{Q-v_2K}{v_1-v_2},\qquad
 k_2=\frac{v_1K-Q}{v_1-v_2}. \label{eq:two}
\end{equation}
Thus two different free-flow lane occupations with the same $K$ but different
$Q$ do not demonstrate information hidden from this observer.
For equal $v_1=v_2=v$, the instantaneous flow becomes $Q=vK$ and no longer
adds an independent equation. Selecting exact equality solely to obtain this
ambiguity would create a symmetry-dependent claim. It is not approved here.
Distinct speeds arbitrarily near equality restore algebraic invertibility
in the noiseless map; numerical conditioning and noisy estimation are separate
questions from exact identity.

\paragraph{Speed has a measurement meaning.}
A density-weighted space mean is redundant with \eqref{eq:moments}:
\begin{equation}
 V_s=Q/K, \qquad K>0.
\end{equation}
In the same equilibrium, point-detector subcase, a vehicle-passage-weighted
arithmetic speed instead supplies
\begin{equation}
 V_t=\frac{\sum_\ell q_\ell v_\ell}{Q}
     =\frac{\sum_\ell v_\ell^2 k_\ell}{Q},\qquad Q>0.
\end{equation}
For three lanes, distinct known speeds and these co-supported observations,
\begin{equation}
 \begin{bmatrix}1&1&1\\v_1&v_2&v_3\\v_1^2&v_2^2&v_3^2\end{bmatrix}
 \begin{bmatrix}k_1\\k_2\\k_3\end{bmatrix}
 =\begin{bmatrix}K\\Q\\QV_t\end{bmatrix},\qquad
 \det A=(v_2-v_1)(v_3-v_1)(v_3-v_2)\ne0.
 \label{eq:three}
\end{equation}
This is an identifiability result for that specific observation map, not a
field estimator or a statement about mismatched station/cell averages.
A road-average speed must not silently be replaced with $Q/K$ if it means
an independent detector statistic. Source [1] distinguishes time/space mean
speeds and calibrates distinct lane speeds. Its detector-derived density may
not be an independent exact measurement of true cell density; the audit must
keep that dependence explicit.

For three lanes with $K,Q$ alone, adding shoulder $k_1$ determines the remaining
pair when $v_2\ne v_3$. This is a \emph{lossless reparameterization} of three
lane densities, not an autonomous two-dimensional group compression.
If $V_t$ is also admitted under \eqref{eq:three}, no additional split information
is required in this free-flow diagnostic. Congested modes and other measurement
maps require a different proof. Snapshot ambiguity in a higher-dimensional
profile still does not establish ambiguity of its full history.

\paragraph{History and memory.}
The required certificate would compare two realizable trajectories with
\begin{equation}
 H(x(t))=H(x'(t))\quad\text{at every declared observed time/cell},
 \qquad u_{\rm past},d,\text{geometry and parameters common}. \label{eq:history}
\end{equation}
No trajectory satisfying \eqref{eq:history} and differing useful metering
responses is constructed in this audit. An unobserved prehistory, changed lane
parameters or unequal speed aggregates cannot be concealed in such a pair.
Likewise, a deterministic hysteresis mode updated from admitted density history
and common initial mode [3] is reconstructible memory, not an extra hidden state.
The field mechanism [4] does not by itself settle this identifiability issue.

\section{Meter permission is not guaranteed physical entry}
Let $u$ be meter-permitted rate; $d_r$ the local ramp sending offer at the merge,
affected by meter releases, ramp propagation and stored traffic; $d_s$ the
shoulder-mainline offer; $r,f_s$ their realized entry rates; $S$ finite
receiving rate. At a common receiving point necessary constraints include
\begin{equation}
 0\le r\le d_r,\qquad 0\le f_s\le d_s,\qquad r+f_s\le S.
 \label{eq:budget}
\end{equation}
An upstream meter constrains its own release $e_r(t)\le u(t)$, not necessarily
instantaneous merge entry $r(t)$: released traffic can arrive later. Only an
explicitly co-located offer gate gives $d_r(t)\le u(t)$. That special mapping is
not adopted here. Account for post-meter in-flight mass if it exists.

At the full multilane interface, every incoming longitudinal/lateral transfer
must share its target receiving budget; every outgoing transfer must also respect
its donor budget. Equation \eqref{eq:budget} isolates a two-input component;
it does not authorize reserving supply twice or prescribe an update order.

With two offers exceeding shared supply, these inequalities alone admit many
allocations. Conservation and a finite capacity do not select the ramp/mainline
service split. A consistent priority/redistribution rule is needed before
there is a unique forward map, let alone its derivative. The ramp-control
source [2] prioritizes ramp flow and excludes mainline spillback onto ramps.
Its assumption cannot replace a finite-access contract for arbitrary demand.
Finite-receiving ramp preference itself remains a valid alternative when its
bounds and domain are made explicit.

The primary node framework [5] describes fixed-priority allocation and unused
supply redistribution. It cautions against setting priorities proportional to
current demand because of nodal invariance. Therefore the neutral proportional
\emph{batch} allocation used in the separate VSL analytical construction cannot
be inherited as a universally sound continuous-flow merge law. No change to that
closed VSL result follows: the traffic semantics and scope differ.

\section{Complete accounting rejects a queue-location-only benefit}
Let $M(t)$ be all uncompleted affected traffic inside the accounting boundary,
including ramp and external waiting, mainline cell masses and residual demand.
For common total arrivals $A(t)$ and departure rate $E(t)$,
\begin{equation}
 \dot M(t)=A(t)-E(t),\qquad
 J=\int_0^{T_c}M(t)\,dt,\qquad M(T_c)=0. \label{eq:account}
\end{equation}
$J$ is total system time, with a common sufficiently long clearance horizon
$T_c$; a delay objective may subtract a policy-independent reference time.
If two policies have common initial total mass, arrivals and full departure
history, they have identical $M(t)$ and $J$. Moving waiting from the ramp to
the mainline alone cannot improve this objective.

On a fixed differentiable branch with control parameter $\theta$, common forcing
and common terminal accounting, the sensitivity is
\begin{equation}
 \frac{dJ}{d\theta}
 =-\int_0^{T_c}(T_c-t)\frac{\partial E(t;\theta)}{\partial\theta}\,dt.
 \label{eq:grad}
\end{equation}
This follows by integrating the sensitivity of \eqref{eq:account} and exchanging
the order of integration. It is not an independent gradient verification.
If a work-conserving fixed-supply bottleneck remains continuously busy with
unchanged aggregate departures for all compared perturbations, the derivative
is zero even if local queues differ. A useful rate decision must change actual
departure timing, future service/spillback, or another explicitly approved
outcome. Merely introducing finite shared receiving does not ensure it.

\section{Disposition and next gate}
The compact family is not refuted. However, autonomous group closure is not
inherited, lane contrast need not be hidden from aggregate flow/speed, and queue
redistribution alone does not establish an all-stream action benefit.
The complete source scheme has unresolved branch/flow-budget transcription
issues, while the ramp extension uses a priority assumption outside the general
finite-access contract. No equation repair or symmetry assumption is adopted.

\textbf{RM9-UG1 asks for a finite-merge allocation contract:} start with
source-backed demand-independent priorities plus unused-supply redistribution
(recommended), use explicitly bounded ramp preference, or adopt a deliberately
stylized proportional-cohort law with a narrower claim. After that decision,
continue a source-to-conservative-formulation audit and test exact information
matching before attempting an action/gradient certificate. Priorities' numerical
values, spatial discretization and lane count are not frozen by this report.
All configurations must preserve the strong observer. No experiment is required
or authorized by this gate.

\begin{thebibliography}{9}
\bibitem{lane} Nagalur Subraveti, Knoop, van Arem.
\emph{First order multi-lane traffic flow model: an incentive based macroscopic
model to represent lane change dynamics}. DOI
\href{https://doi.org/10.1080/21680566.2019.1700846}{10.1080/21680566.2019.1700846}.
FB0030 final PDF, printed pp. 1760--1766 and 1768--1769; equations 1--22,
Tables 1--2. Publication/repository year labels differ (2019/2020).
\bibitem{ramp} Nagalur Subraveti, Knoop, van Arem.
\emph{Lane Change Control Combined with Ramp Metering}. 2021.
\href{https://doi.org/10.1155/2021/8638523}{10.1155/2021/8638523}.
RM0004, Sections 2.1--3. Final repository PDF inspected.
\bibitem{memory} Cao, Nilsson, Coogan.
\emph{On the Impact of the Capacity Drop Phenomenon for Freeway Traffic Flow
Control}. 2021. \href{https://arxiv.org/abs/2106.10347}{arXiv:2106.10347}.
RM0003, Section II.B, equation 3; archived PDF and public methods.
\bibitem{field} Cassidy, Rudjanakanoknad.
\emph{Increasing the capacity of an isolated merge by metering its on-ramp}.
RM0014, public 2004 report, Sections 3--5; distinct from 2005 journal version.
\href{https://escholarship.org/uc/item/30s46939}{Primary report landing page}.
\bibitem{node} Wright, Gomes, Horowitz, Kurzhanskiy.
\emph{On node models for high-dimensional road networks}.
Transportation Research B 105 (2017), 212--234.
RM0018, Section 3.1, Remark 3.1 and Section 4.2 priority examples.
\href{https://horowitz.me.berkeley.edu/Publications_files/All_papers_numbered/Wright_node_models_transportation_networks_Transportation_research_B_2017.pdf}{Author-hosted published PDF}.
\end{thebibliography}
\end{document}
